\subsection{Polynomial functions on an infinite integral domain}

PROPOSITION 9. --- Assume that A is an integral domain. Let I be a set, \((\mathrm{H}_{i})_{i\in I}\) a family of infinite subsets of A and \(H=\prod_{i\in I}H_{i}\subset A^{I}\) . If f is a non-zero

element of \(\mathbf{A}[(\mathbf{X}_i)_{i\in I}]\) , and \(\mathrm{H}_f\) the set of all \(\mathbf{x}\in \mathrm{H}\) such that \(f(\mathbf{x})\neq 0\) , then \(\mathrm{H}\) and \(\mathrm{H}_f\) are equipotent.

\begin{enumerate}
\def\labelenumi{\alph{enumi})}
\tightlist
\item
  First suppose that I is finite and put \(n = \text{Card I}\) . The proposition is clear for \(n = 0\) ; we shall prove it by induction on \(n\) . Choose an element \(i_0\) of I and put \(J = I - \{i_0\}\) , \(B = A[(X_i)_{i \in J}]\) . Since \(f \neq 0\) , we can write \(f = \sum_{k=0}^{m} g_k X_{i_0}^k\) , where \(g_0, \ldots, g_m \in \mathbf{B}\) and \(g_m \neq 0\) . By the induction hypothesis the set \(\mathbf{K}\) of all \(\mathbf{x} \in \prod_{i \in J} \mathrm{H}_i\) such that \(g_m(\mathbf{x}) \neq 0\) is equipotent with \(\prod_{i \in J} \mathrm{H}_i\) . For \(\mathbf{x} \in \mathbf{K}\) the polynomial
\end{enumerate}

\[
h \left(\mathrm{X} _ {i _ {0}}\right) = \sum_ {k = 0} ^ {m} g _ {k} (\mathbf {x}) \mathrm{X} _ {i _ {0}} ^ {k} \in \mathrm{A} \left[ \mathrm{X} _ {i _ {0}} \right]
\]

is non-zero. By Th. 1 (IV, p.~16) the set of \(\alpha \in \mathrm{H}_{i_0}\) such that \(h(\alpha) \neq 0\) is equipotent with \(\mathrm{H}_{i_0}\) whence

\[
\operatorname{Card} \mathrm{H} \geqslant \operatorname{Card} \mathrm{H} _ {f} \geqslant (\operatorname{Card} \mathrm{K}). (\operatorname{Card} \mathrm{H} _ {i _ {0}}) = \operatorname{Card} \mathrm{H},
\]

and so Card \(\mathbf{H} = \mathbf{Card}\mathbf{H}_f\)

\begin{enumerate}
\def\labelenumi{\alph{enumi})}
\setcounter{enumi}{1}
\tightlist
\item
  In the general case there is a finite subset \(\mathbf{I}'\) of \(\mathbf{I}\) such that \(f \in \mathbf{A}[(\mathbf{X}_i)_{i \in \mathbf{I}'}]\) . Let \(\mathbf{H}_f'\) be the set of all \(\mathbf{x} \in \prod_{i \in \mathbf{I}'} \mathbf{H}_i\) such that \(f(\mathbf{x}) \neq 0\) . Then \(\mathbf{H}_f = \mathbf{H}_f' \times \left( \prod_{i \in \mathbf{I} - \mathbf{I}'} \mathbf{H}_i \right)\) , and it suffices to apply the first part of the proof to \(\mathbf{H}_f'\) .
\end{enumerate}

COROLLARY 1. --- We keep the hypothesis and notation of Prop. 9. If I is non-empty, then \(H_{f}\) is infinite.

COROLLARY 2. --- Suppose that A is an infinite integral domain or that A is an algebra over an infinite field. For every \(f \in \mathbf{A}[(\mathbf{X}_i)_{i \in I}]\) , let \(\tilde{f}: \mathbf{A}^{\mathrm{I}} \to \mathbf{A}\) be the polynomial function defined by \(f\) (IV, p.~4). Then the mapping \(f \mapsto \tilde{f}\) is injective.

When \(\mathbf{A}\) is an infinite integral domain, the corollary follows at once from Prop. 9. Suppose that \(\mathbf{A}\) is an algebra over an infinite field \(k\) . Let \(f = \sum_{\nu \in \mathbb{N}^{(I)}} \alpha_{\nu} X^{\nu}\)

be a non-zero element of \(\mathbf{A}[(\mathbf{X}_i)_{i\in I}]\) ; then there exists \(\nu_0\in \mathbf{N}^{(I)}\) such that \(\alpha_{\nu_0}\neq 0\) , and a \(k\) -linear form \(\varphi\) on \(\mathbf{A}\) such that \(\varphi (\alpha_{\nu_0})\neq 0\) . Let \(g = \sum_{\nu \in \mathbf{N}^{(I)}}\varphi (\alpha_{\nu})\mathbf{X}^{\nu}\in k[(\mathbf{X}_i)_{i\in I}]\) ; we have \(g\neq 0\) , hence there exists \(\mathbf{x}\in k^{\mathrm{I}}\) such

\[
g (\mathbf {x}) \neq 0
\]

\[
\varphi (f (\mathbf {x})) = g (\mathbf {x}) \neq 0
\]

\[
f (\mathbf {x}) \neq 0
\]

When A is an infinite integral domain or when A is an algebra over an infinite field, we shall usually identify f with \(\tilde{f}\) .

Suppose that A is finite and let \(f(X) = \prod_{\alpha \in A} (X - \alpha)\) , then \(f \neq 0\) , but \(\tilde{f} = 0\) . For other examples see IV, p.~88, exercises 7 and 8.

THEOREM 2 (Principle of extension of algebraic identities). --- Suppose that A is an infinite integral domain. Let \(g_1, \ldots, g_m\) , \(f\) be elements of \(\mathbf{A}[(\mathbf{X}_i)_{i \in I}]\) and assume the following hypotheses:

\begin{enumerate}
\def\labelenumi{\alph{enumi})}
\item
  \(g_{1} \neq 0, \ldots, g_{m} \neq 0\) ;
\item
  for all \(\mathbf{x} \in \mathbf{A}^{I}\) such that \(g_1(\mathbf{x}) \neq 0, \ldots, g_m(\mathbf{x}) \neq 0\) , we have \(f(\mathbf{x}) = 0\) . Then \(f = 0\) .
\end{enumerate}

For if \(f \neq 0\) , we have \(fg_1 \ldots g_m \neq 0\) (IV, p.~9, Prop. 8), hence there exists \(\mathbf{x} \in \mathbf{A}^{I}\) such that \(f(\mathbf{x})g_1(\mathbf{x}) \ldots g_m(\mathbf{x}) \neq 0\) (IV, p.~18, Cor. 2), which contradicts the hypothesis.

Scholium. --- Let A be an integral domain and \(f \in \mathbf{A}[(\mathbf{X}_i)_{i \in I}]\) . Th. 2 provides a convenient means of proving that f = 0. It suffices to consider an infinite integral domain E containing A as subring; if we can show that \(f((x_i)) = 0\) for all \((x_i) \in \mathrm{E}^{\mathrm{I}}\) (or even for those \((x_i) \in \mathrm{E}^{\mathrm{I}}\) at which a finite number of given non-zero polynomials do not vanish) then it follows that f = 0. If A itself is not infinite, we can for example take E to be the ring A{[}X{]} or its field of fractions.

Once we have proved the relation \(f = 0\) , we can clearly deduce that \(f((y_i)) = 0\) for all \((y_i) \in \mathbf{F}^{I}\) where \(\mathbf{F}\) is any unital associative and commutative \(\mathbf{A}\) -algebra whatsoever; in particular \(\mathbf{F}\) may be finite or non-integral.

In other words, the proof of the identity \(f((x_{i})) = 0\) when the \(x_{i}\) run over an infinite integral domain containing A as subring (with the possible restriction that \(g_{k}((x_{i})) \neq 0\) for \(1 \leqslant k \leqslant m\) , where the \(g_{k}\) are non-zero polynomials) implies the same identity when the \(x_{i}\) run over any unital associative and commutative A-algebra.

In particular let \(f \in \mathbf{Z}[(\mathbf{X}_i)_{i \in I}]\) . If \(f((x_i)) = 0\) when the \(x_i\) run over \(\mathbf{Z}\) (with the possible restriction that \(g_k((x_i)) \neq 0\) for \(1 \leqslant k \leqslant m\) , where the \(g_k\) are non-zero elements of \(\mathbf{Z}[(\mathbf{X}_i)]\) ), then we have the same identity when the \(x_i\) run over an arbitrary commutative ring.

